

Throughout our proofs, we will be forming and manipulating higher-order
Taylor expansions of functions of $\theta$.  The arrays of derivatives
in these expressions will always be summed against a direction, typically
$\tau = \sqrt{N}(\theta - \thetahat)$.  In order to avoid cumbersome
summation notation for all such terms, we will implicitly sum derivative arrays
against $\tau$ or $\theta$.  or example, for
a function $\psi(\theta, \x_n) \in \rdom{D_\psi}$, we will write
%
\begin{align*}
\psigrad{3}(\thetahat, \x_n)\tau^3 =
\sum_{a=1}^D \sum_{b=1}^D \sum_{c=1}^D
    \fracat{\partial^3 \psi(\theta, \x_n)}
           {\partial\theta_{a} \partial\theta_{b} \partial\theta_{c}}
           {\thetahat} \tau_{a} \tau_{b} \tau_{c}.
\end{align*}
%
Occasionally we will need to sum several derivative arrays against $\tau$,
particularly in \secref{region_r1}.  Since the directions $\tau$ will always be
the same, we will combine the $\tau$ in such expressions with no ambiguity. For
example, for the term involving $\likhatk{3}(\thetahat)$ in the statement of
\thmref{bayes_clt_main}, we will write
%
\begin{align*}
%
\psigrad{1}(\thetahat, \x) \likhatk{3}(\thetahat) \tau^4 =
\sum_{a=1}^{D} \sum_{b=1}^{D} \sum_{c=1}^{D} \sum_{d=1}^{D}
    \psigrad{1:a}(\thetahat, \x) \likhatk{3:b,c,d}(\thetahat)
    \tau_a \tau_b \tau_c \tau_d.
%
\end{align*}
%

When necessary, higher
derivatives will be indexed in the parentheses.  For example,
$\psigrad{2:i,j}(\theta, \x_n)$ will denote the length--$D_\psi$ vector
%
\begin{align*}
%
\fracat{\partial^2 \psi(\theta, \x_n)}
    {\partial\theta_i \partial\theta_j}{\theta, \x_n}.
%
\end{align*}
%
%\end{defn}

We denote by $\ind{A}$ the indicator function taking value $1$ when the
event $A$ occurs and $0$ otherwise.

% Order notation
We will use order notation to describe the behavior of quantities as $N
\rightarrow \infty$.  For example, we will use $z = O(\phi(N))$ to mean there
exists an $N^*$ and an $A > 0$ such that $N > N^* \Rightarrow z < A \phi(N)$. In
our notation, $A$ and $N^*$ in this definition will never depend on the data,
$\xvec$, though they may depend on unobserved population quantities.
%
Dependence on the data $\xvec$ in order notation will always be represented with
$O_p$ notation.  For example, $z = O_p(\phi(N))$ will mean that there exists an
$A > 0$ such that for any $\epsilon > 0$, there is a corresponding $N^*$ and
such that
%
\begin{align*}
%
N > N^* \quad\Rightarrow\quad
    \expect{\xfdist}{\ind{z < A \phi(N)}} > 1 - \epsilon.
%
\end{align*}
%
Here, $A$ does not depend on $\epsilon$.  The meaning is that the event $z =
O(\phi(N))$ occurs with arbitrarily high probability as $N$ goes to infinity.
The notation $\tilde{O}$ and $\tilde{O}_p$ will have the same meaning as $O$ and
$O_p$ notation respectively, but with factors of $\log N$ possibly ignored.
